\documentclass[tikz,border=4pt]{standalone}
\usepackage{ctex}
\usepackage{amsmath,amssymb}
\usetikzlibrary{calc,arrows.meta,decorations.markings}
\begin{document}
\begin{tikzpicture}[
  x={(1cm,0cm)}, y={(0cm,1cm)}, z={(-0.48cm,-0.38cm)},
  line cap=round, line join=round,
  >=Stealth,
  decoration={markings, mark=at position 0.5 with {\arrow{Stealth}}}
]
% 坐标轴
\draw[->] (0,0,0) -- (4.5,0,0) node[right] {$x$};
\draw[->] (0,0,0) -- (0,4,0) node[above] {$y$};
\draw[->] (0,0,0) -- (0,0,3.5) node[above left] {$z$};

% 曲面 S（截取的曲面，简化为矩形参数曲面）
\fill[blue!20, opacity=0.7]
  (0.5,0.5,1.5) -- (3.5,0.5,2.0) -- (3.5,3.0,2.5) -- (0.5,3.0,2.0) -- cycle;
\draw[blue!60, thick]
  (0.5,0.5,1.5) -- (3.5,0.5,2.0) -- (3.5,3.0,2.5) -- (0.5,3.0,2.0) -- cycle;
\node[blue!70!black, font=\small] at (2.0,1.8,2.2) {$S$};

% 边界曲线 partial S（逆时针，对应 n 朝上）
\draw[red, very thick, postaction={decorate}]
  (0.5,0.5,1.5) -- (3.5,0.5,2.0);
\draw[red, very thick, postaction={decorate}]
  (3.5,0.5,2.0) -- (3.5,3.0,2.5);
\draw[red, very thick, postaction={decorate}]
  (3.5,3.0,2.5) -- (0.5,3.0,2.0);
\draw[red, very thick, postaction={decorate}]
  (0.5,3.0,2.0) -- (0.5,0.5,1.5);
\node[red, font=\small] at (4.0,0.3,2.1) {$\partial S$};

% 法向量 n
\draw[->, orange, ultra thick]
  (2.0,1.75,2.05) -- ++(0,0,1.2)
  node[right, font=\small] {$\mathbf{n}$};
\node[orange!80!black, font=\small] at (2.8,0.5,3.2) {右手定则};

% 旋度向量（curl F）示意
\draw[->, purple!70, thick]
  (1.5,2.5,2.2) -- ++(-0.2,0,0.8)
  node[above, font=\small] {$\nabla\times\mathbf{F}$};

% Stokes 公式标注
\node[font=\small, draw=gray!50, fill=gray!10, rounded corners=2pt, align=center]
  at (-0.8,4,3.5)
  {\textbf{Stokes 定理}\\
   $\oint_{\partial S}\mathbf{F}\cdot d\mathbf{r} = \iint_S(\nabla\times\mathbf{F})\cdot d\mathbf{S}$\\
   边界方向与法向满足右手定则};
\end{tikzpicture}
\end{document}
